\section{Introduction}
\label{sec:intro}

Compactness has been a criterion in redistricting law since at least the
late 19th century.  Courts, legislatures, and redistricting commissions
routinely cite compactness as a marker of a non-gerrymandered map.  The
intuition: a compact district is one that ``looks like'' a geographic region,
not a political construction.  A district that winds around cities to include
partisan voters while excluding others is, by inspection, not compact.

But compactness is more ambiguous than this intuition suggests.  First,
there are many measures of compactness---Polsby-Popper \citep{polsby1991third},
Reock \citep{reock1961note}, convex hull ratio, length-width ratio---and they
do not always agree.  Second, compactness is not neutral in sorted states:
more compact plans tend to produce more Republican seats, because Democratic
voters are geographically concentrated in cities \citep{rodden2019why}.
Third, an extremely compact plan can be a form of ``packing'': concentrating
Democratic voters in a small number of very compact urban districts to
maximise the number of Republican districts.

The \ar{} algorithm \citep{dellaLibera2026b11} optimises the METIS edge-cut
objective, which is correlated with but not identical to Polsby-Popper
compactness.  The algorithm does not directly maximise compactness; it
minimises the number of cross-district census-tract boundaries.  The resulting
plans are compact but not maximally compact.

\paragraph{Main questions.}
\begin{enumerate}
  \item Where does the \ar{} plan fall in the \recom{} ensemble's compactness
    distribution?
  \item Is the \ar{} plan an outlier in the legal sense?
  \item How does compactness correlate with partisan outcome within the ensemble,
    and what does this imply for the \ar{} plan's legal defensibility?
\end{enumerate}

\paragraph{Main findings.}
\begin{enumerate}
  \item The \ar{} plan falls at the $61$st--$72$nd percentile of mean
    Polsby-Popper score across six states.  It is more compact than average
    but not maximally compact.
  \item The \ar{} plan is not an outlier in the legal sense: it is not at
    the 95th percentile or above on any compactness measure.
  \item Within the \recom{} ensemble, higher compactness tends to correlate
    with more Republican seats in sorted states (the Rodden effect).  The
    \ar{} plan's compactness is therefore consistent with---but does not
    cause---its moderate Republican lean in sorted states.
\end{enumerate}

\paragraph{Organisation.}
Section~\ref{sec:distribution} characterises the \recom{} compactness
distribution.  Section~\ref{sec:pp-percentile} reports the \ar{} plan's
PP percentile.  Section~\ref{sec:outlier} addresses the outlier question.
Section~\ref{sec:correlation} examines the compactness-partisan correlation.
Section~\ref{sec:legal} develops the legal defensibility argument.
Section~\ref{sec:conclusion} concludes.
